% =====================================================================
% Web-search retrieval arm vs data-lake arms, LakeQA subset20b.
%
% Generated from experiments/2026-08-31-web-arm-subset20b/
%   results/            raw exact_match
%   results_semantic/   semantic_match (judge: gpt-5.4, effort medium)
% Reproduce with summarize.py / per_task.csv.
%
% Requires: booktabs.
% =====================================================================

\subsection{Retrieval source: open web versus the data lake}
\label{sec:web-arm}

To test whether the retrieval axis is specific to the data lake or reflects a
more general property of the agent, we replaced lake search with open-web
search. The arm uses the Parallel Search API and, under \texttt{-{}-no-s3},
loses every lake-backed file tool: the agent retains only \texttt{download},
which fetches an arbitrary \texttt{http(s)} URL into its sandbox, and
\texttt{execute\_code}. All four arms hold the remaining axes fixed at
\texttt{search\_results=naive}, \texttt{profile=standard} and
\texttt{computation\_tool=standard}, since the ideal computation tools resolve
against authored data-lake records and would make a web run's outcome
independent of what it retrieved.

Table~\ref{tab:web-arm} reports the result. The web arm reaches 35.0\%
semantic match against 45--50\% for the three lake arms, and the three lake
conditions are separated by a single task. The interesting quantity is not the
gap in accuracy but the gap in retrieval: the ideal arm reaches 78.1\% of the
gold sources using 22.2 cycles, while the standard arm reaches 67.9\% using
29.2 cycles and 56\% more money. Oracle retrieval is markedly more efficient
and yet converts to the same number of correct answers, which locates the
bottleneck downstream of retrieval on this split.

Two properties of the web arm complicate its interpretation and we report them
rather than correct for them. First, \texttt{download} accepts any URL, and the
model routinely constructs one from pretrained knowledge rather than following
a search result --- live Socrata queries against
\texttt{data.cityofnewyork.us}, \texttt{data.va.gov} export endpoints, and in
one case a timestamped \texttt{web.archive.org} URL. The arm therefore measures
an agent with open web access and code execution, not web-search retrieval, and
its data path may be \emph{stronger} than the lake arms', since it aggregates
live authoritative sources rather than a fixed snapshot. Second, that liveness
produces false negatives: two failures are arithmetically correct against
current data but disagree with gold answers derived from the snapshot
(2{,}700{,}968 against a gold 2{,}705{,}988; 2923.59 against 2941).

A third of all task-runs (33 of 80) reached the 30-call tool budget, and every
arm scored roughly two to three times better below the cap than at it, so the
absolute numbers here should be read as depressed by the budget rather than as
ceilings.

% ---------------------------------------------------------------------
\begin{table}[h]
  \centering
  \caption{Retrieval-source comparison on LakeQA \textsc{subset20b}
  (\texttt{gpt-5-mini}, 20 tasks per arm, zero errored runs). SM is semantic
  match (judge \texttt{gpt-5.4}); EM is lexical exact match. \emph{Blank} counts
  answers the judge classified as \texttt{answer\_unknown\_blank}, i.e.\ honest
  abstentions rather than wrong answers. \emph{Src} is the fraction of required
  gold sources the agent read; the web arm touches no lake sources by
  construction. \emph{Cap} counts tasks that reached the 30-call tool budget.}
  \label{tab:web-arm}
  \scriptsize
  \setlength{\tabcolsep}{4pt}
  \renewcommand{\arraystretch}{0.95}
  \resizebox{\columnwidth}{!}{%
  \begin{tabular}{lrrrrrrrr}
    \toprule
    Arm & SM (\%) & EM (\%) & Blank & Cycles & Tools & Cap & Cost (\$) & Src (\%) \\
    \midrule
    Ideal    & \textbf{50.0} & 45.0 & 0 & 22.2 & 21.8 & 5  & 0.60 & \textbf{78.1} \\
    Standard & \textbf{50.0} & 45.0 & 0 & 29.2 & 28.2 & 11 & 0.94 & 67.9 \\
    Naive    & 45.0 & 45.0 & 0 & 27.8 & 26.8 & 9  & 0.83 & 70.1 \\
    Web      & 35.0 & 30.0 & 1 & 26.2 & 25.2 & 8  & 1.00 & --- \\
    \bottomrule
  \end{tabular}}
\end{table}

% ---------------------------------------------------------------------
\paragraph{Semantic match versus exact match.}
The two metrics differ by exactly one task, and that task is instructive. For
the question whose gold answer is \emph{``the Cut Bank Penguin''}, three arms
answered \texttt{[Cut Bank Penguin]} and the fourth answered
\texttt{[World's largest statue\ldots]}. All four are correct; all four score
zero under exact match, because the lexical metric penalises a dropped article.
On LakeQA most gold answers are a single token, where token-overlap F1 is
identical to exact match by construction, and where it does differ it largely
measures whether the model reproduced determiners. We therefore report semantic
match throughout and do not report F1.
